\section{Compatible representation families and canonical denominator data}
\label{sec:denominator-family-v051}

\status{DFD1--DFD9 construct the actual continuous integral complements,
the original \texttt{Pt\allowbreak Fam}, and the whole original
\texttt{exists\_\allowbreak hub\allowbreak Data} output. DFD10--DFD14 supply the good-row
selection, the exact twist-killing implication, the finite-coefficient
reduction of the retained integral crystalline theorem, and uniform
parameter and small-prime bounds. DFD15 assembles the original family
setup conditional on two explicitly stated remaining suppliers: the
coherent Wach comparison and the Selmer signature comparison. The full
\texttt{exists\_\allowbreak setup\_\allowbreak large} producer and \texttt{Concl\allowbreak F} remain open in
the human reconstruction.}

Fix the original datum $D:\texttt{D7}$. Its host is the same CM field
$E$, its absolute Galois group is $G_E$, and its rank parameters are
$N=n_V=r+t$. At a location $i$, write $O_i$ for its coefficient DVR,
$K_i=\operatorname{Frac}(O_i)$, $\ell_i=D.\mathrm{nd}.\ell(i)$,
$X_i$ for the point set and $Z_i\subseteq X_i$ for the specified
automorphic points. The characters $\chi_{i,c}:G_E\to O_i^\times$,
for $1\le c\le t$, and both coefficient embeddings are those of $D$.
An integral representation will mean a homomorphism into a general
linear group whose matrix entries are continuous. Its rational
semisimplicity is always stated separately.

The original \texttt{Cheb} premise is retained wherever used. It
asserts the degree-one, unramified Frobenius approximation with finite
avoidance in every finite normal quotient of $G_E$. The present
argument does not replace that premise by a density statement for an
unspecified set of places. Likewise \texttt{adm2},
\texttt{Conclusion41}, \texttt{Facts\allowbreak F}, and the local comparison data
below retain their original source meanings.

\subsection{From trace identities to actual integral complements}

% BEGIN DFD1
\paragraph{DFD1: a character summand descends to the original field.}
\label{claim:DFD1}
Let $K\hookrightarrow L$ be an embedding of fields with $L$ of
characteristic zero, $G$ an arbitrary group, and $V$ a finite-dimensional
semisimple $K[G]$-module. Suppose $T:G\to L$ is the trace of an actual
finite-dimensional $L$-representation and
\[
 \operatorname{tr}(g\mid V)_L=T(g)+\sum_{j=1}^{t}\chi_j(g),
 \qquad \chi_j:G\longrightarrow K^\times.                 \tag{DFD.1}
\]
Then there are independent vectors $v_1,\ldots,v_t\in V$ with
$gv_j=\chi_j(g)v_j$, and a $G$-stable complement $W$ such that
\[
 V=W\oplus\bigoplus_{j=1}^{t}Kv_j.                       \tag{DFD.2}
\]
The characters may repeat; no continuity of $K\hookrightarrow L$ is
required for this algebraic assertion.

Here is the trace argument, including the descent. Extend scalars to
an algebraic closure of $L$. Replace each representation by its
semisimplification when computing its trace: a matrix preserving a
finite composition series is block triangular, so its trace is the
sum of its traces on the composition factors. For the finite direct
sum of all the resulting representations, the image of the group
algebra is a finite-dimensional semisimple algebra. Indeed it acts
faithfully on a semisimple module, so its radical acts by zero and
is zero. Over the algebraically closed field it is a product of
matrix algebras. The trace on the central idempotent of one factor
is the dimension of its simple module times its multiplicity.
The equality (DFD.1), first extended linearly from group elements to
the group algebra and then evaluated on a lift of each idempotent,
therefore identifies the multiplicities. Characteristic zero makes
the positive dimension invertible. Every occurrence of a $\chi_j$
on the right consequently occurs on the left.

Scalar extension preserves semisimplicity here: the image algebra
on $V$ is semisimple; its simple factors have separable centers over
$K$, since $K$ has characteristic zero, and their scalar extensions
are semisimple. In particular the occurrence just found supplies
a nonzero $\chi_1$-eigenvector in the scalar extension of $V$.
The simultaneous equations
\[
 (\rho(g)-\chi_1(g))v=0\quad(g\in G)
\]
have coefficients in $K$. Their row space is finite-dimensional,
so finitely many rows already span all the rows. The rank of this
finite matrix is unchanged by field extension. Its kernel is
therefore nonzero over $K$ as well. Choose $v_1$ there.

Semisimplicity gives a stable complement to $Kv_1$. In a basis
adapted to this decomposition the trace is
$\chi_1+\operatorname{tr}(W_1)$. Subtract $\chi_1$ from (DFD.1)
and repeat on $W_1$. At the $j$th stage the new vector belongs to
the complement of the previously chosen lines, so independence is
automatic, including for repeated characters. After $t$ stages
the remaining stable summand is the required $W$. This proves the
full eigenvector and rational splitting assertions used by the
source's \texttt{exists\_\allowbreak eigenvectors} and \texttt{exists\_\allowbreak complement}.
% END DFD1

% BEGIN DFD2
\paragraph{DFD2: a saturated lattice gives a continuous integral complement.}
\label{claim:DFD2}
Let $O$ be a topological principal ideal domain with fraction field
$K$, and let $R:G\to\operatorname{GL}_N(O)$ have continuous matrix
entries. Suppose $R_K$ is semisimple and its trace has the form
(DFD.1), with $N=a+t$ and the characters valued in $O^\times$.
There is a continuous $\rho:G\to\operatorname{GL}_a(O)$ and an
intertwining $K$-linear isomorphism
\[
 K^N\simeq K^a\oplus K^t,
 \qquad R_K\simeq\rho_K\oplus
      \operatorname{diag}(\chi_1,\ldots,\chi_t),           \tag{DFD.3}
\]
whose block representation is semisimple.

Take the stable summand $W\subset K^N$ from DFD1 and set
$\Lambda=W\cap O^N$. It is finitely generated because $O$ is
Noetherian. It is saturated: if $0\ne c\in O$, $v\in O^N$ and
$cv\in\Lambda$, then $v\in W$ because $W$ is a $K$-space, hence
$v\in\Lambda$. Thus $O^N/\Lambda$ is torsion-free and finitely
generated, and so free over the PID. The exact sequence splits
as $O$-modules. Choose a basis of $\Lambda$ and an $O$-linear
retraction $q:O^N\to O^a$ of its inclusion
$b:O^a\to O^N$. Its rank is $a$: clear the finitely many
denominators of a $K$-basis of $W$ to see that $K\Lambda=W$.

Both $R(g)$ and $R(g^{-1})$ preserve $\Lambda$. Define
\[
             \rho(g)=qR(g)b.                             \tag{DFD.4}
\]
On $\Lambda$, $bq$ is the identity. This proves
$\rho(gh)=\rho(g)\rho(h)$ and
$\rho(g^{-1})=\rho(g)^{-1}$. Every matrix entry of (DFD.4) is a
finite $O$-linear combination of entries of $R(g)$, so it is
continuous. This argument uses continuity only inside $O$;
there is no assumption that a chosen rational basis change is
integral or continuous in an unspecified coefficient topology.

The map $b_K$ identifies $K^a$ with $W$. Adjoin the character
vectors of DFD1 to obtain (DFD.3). Semisimplicity transfers along
this isomorphism from $R_K$. These are exactly the three fields
\texttt{rho}, \texttt{ss}, and \texttt{cont} of the original
\texttt{Is\allowbreak Complement}. The original lattice $O^N$ need not be
the direct sum of the character lattices and $\Lambda$; the
asserted block decomposition is rational, as in the source.
% END DFD2

% BEGIN DFD3
\paragraph{DFD3: equality on the permitted Frobenius elements is enough.}
\label{claim:DFD3}
Assume the original \texttt{Cheb}. If two continuous maps
$f,f':G_E\to Y$, with $Y$ Hausdorff, agree on every Frobenius lift
at every split place outside a finite exceptional set, then $f=f'$.
Here ``split'' is the literal host condition
$\texttt{Is\allowbreak Split\allowbreak Pl}$, not a new density convention.

First a place of $E$ unramified over $\mathbb Q$ and of residue
degree one is an allowed split place. The host extension
$E/\mathbb Q$ is Galois. Its decomposition group at that place
has order $ef=1$, so it is trivial. If complex conjugation fixed
the place, it would belong to this group. But the host's
conjugation sends its nonzero square root of $-d$ to its negative,
so it is not the identity. This proves the required split
condition.

If $f(g)\ne f'(g)$, continuity and the Hausdorff property give an
open neighborhood $U$ of $g$ throughout which they differ.
The Krull topology supplies an open normal subgroup $H$ with
$Hg\subseteq U$. Apply \texttt{Cheb} to $H$, $g$, and the finite
exceptional set. It produces an unramified degree-one place and
a Frobenius lift in $Hg$. The preceding paragraph makes that
place split, so equality there contradicts its membership in
$U$. This proves the source's full
\texttt{AE.eq\_\allowbreak of\_\allowbreak continuous} assertion.

For the applications below enlarge the finite exceptional set
by the places over $\ell_i$, the prescribed sets $S_i,Q_i$, and
the finitely many ramification sets of the automorphic parts.
There are finitely many places over $\ell_i$ because they divide
the nonzero ideal $(\ell_i)$ in $\mathcal O_E$. Thus all such
enlargements are legitimate. Equality is extended in
$\overline{\mathbb Q}_{\ell_i}$, with its given Hausdorff
topology. The complex embedding is used only to check the
Frobenius values and is then canceled by injectivity; it is
never assumed continuous.
% END DFD3

% BEGIN DFD4
\paragraph{DFD4: the cuspidal traces supplied by the original admissibility.}
\label{claim:DFD4}
The original $\texttt{adm2}(D)$ implies
$\texttt{Galois\allowbreak Reps\allowbreak Cusp}(D)$ at every location, with the
original embeddings and normalizations. In particular, for a
regular algebraic conjugate-self-dual cuspidal part $\pi$ of
rank $a\ge1$ there is a continuous actual character $T_\pi$
whose cofinite split Frobenius values satisfy
\[
 \iota_{\ell,\mathbb C}T_\pi(\operatorname{Frob}_w)
       =s_\pi(w)q_w^{(1-a)/2},                            \tag{DFD.5}
\]
where $s_\pi(w)$ is the sum of its Satake parameters and
$q_w=\mathrm Nw$.

The first field of \texttt{adm2} supplies the continuous
semisimple Galois representation with its local
Weil--Deligne compatibility away from $\ell_i$. Remove the
finitely many places over $\ell_i$ and the cofinite automorphic
ramification exceptions. At any remaining place the supplied
compatibility identifies the Frobenius characteristic
polynomial with the normalized unramified Satake polynomial.
The coefficient of degree $a-1$ in
$\prod_j(X-\alpha_j)$ is $-\sum_j\alpha_j$, while that of a
matrix characteristic polynomial is minus its trace. Comparing
these coefficients gives (DFD.5). The trace is continuous since
it is a finite sum of diagonal matrix entries; it is an actual
character because it is the trace of this same representation.
Continuity and semisimplicity are unchanged. This proves all
fields of \texttt{Galois\allowbreak Reps\allowbreak Cusp}, and also the exact
\texttt{has\allowbreak Trace\_\allowbreak cusp} consequence. It assumes the Galois
representation already included in \texttt{adm2}; it does not
assert a new automorphy theorem.
% END DFD4

% BEGIN DFD5
\paragraph{DFD5: integral Tate twists in both parity cases.}
\label{claim:DFD5}
Assume \texttt{Galois\allowbreak Reps\allowbreak Cusp}. Let $(a,\pi,d)$ be a part in
the actual isobaric parameter list of a point of $Z_i$. The
normalized trace of its Speh block,
\[
       \sum_{k=0}^{d-1}s_\pi(w)q_w^{(d-N)/2-k},           \tag{DFD.6}
\]
is the cofinite split Frobenius trace of an actual continuous
representation over $\overline{\mathbb Q}_{\ell_i}$.

The infinitesimal parameter of the whole list has distinct
entries in $\mathbb Z+(N-1)/2$. The $k=0$ submultiset of the
chosen part is its infinitesimal parameter translated by
$(d-1)/2$. As a submultiset of the whole parameter it too has
distinct entries. Therefore the infinitesimal weights of
$\pi$ lie in $\mathbb Z+(N-d)/2$ and are distinct. The
distinguished character $\mu=\texttt{mu\allowbreak Eta}(i)$ has weights in
$\mathbb Z+(N-t)/2$. Since $r=k+m-1$ is odd, write
\[
                    N=t+2u+1.                           \tag{DFD.7}
\]
The fixed character $\chi_1$ has Frobenius value
$s_\mu(w)q_w^{(t-N)/2}$; the cyclotomic character has value
$q_w$, and $\chi_{c+1}=\chi_1\operatorname{cyc}^{-c}$.

If $a+d-N=2e+1$, then
$(N-d)/2=(a-1)/2-e$, so $\pi$ is regular algebraic in the
normalization of DFD4. Twist its Galois representation by the
integer Tate power $\operatorname{cyc}^{e-k}$. Its trace has
exponent
\[
                 (1-a)/2+e-k=(d-N)/2-k.
\]
If instead $a+d-N=2e$, twist the automorphic part first by
$\mu$. Its infinitesimal weights differ from $(a-1)/2$ by
integers, since the sum of the two shifts differs by
$u+1-e$. Translation preserves distinctness, and the original
conjugate-self-duality law preserves that condition. Hence
$\pi\otimes\mu$ is an allowed input to DFD4. Twist its
Galois representation by
\[
                \chi_1^{-1}\operatorname{cyc}^{e-u-1-k}.
\]
The factor $s_\mu(w)$ cancels and the exponent becomes
\[
 (1-a)/2+(N-t)/2+e-u-1-k
          =e-a/2-k=(d-N)/2-k.                            \tag{DFD.8}
\]
Every Tate exponent used in both cases is an integer.
The nonzero complex Satake parameter of the character makes
the cancellation legitimate. Powers of the positive real
$q_w$ have the usual additive exponent law, agreeing with the
source complex-power convention.

Finite direct sums over $k$ give (DFD.6), and finite direct
sums over parts give sums of these characters. Tensoring by a
continuous one-dimensional character preserves continuity.
All finite exceptional sets may be united as in DFD3. This is
the full \texttt{has\allowbreak Trace\_\allowbreak part} assertion, including the even
parity case; it requires no half-integral Tate character.
% END DFD5

% BEGIN DFD6
\paragraph{DFD6: the actual block complement at every specified point.}
\label{claim:DFD6}
Under \texttt{Cheb} and \texttt{Galois\allowbreak Reps\allowbreak Cusp}, the original
\texttt{Complement\_\allowbreak at} holds. Namely, for every $i$, every
$\pi\in Z_i$, and every supplied Hecke Galois representation
$R_\pi:G_E\to\operatorname{GL}_N(O_i)$, there is an actual
$\rho_\pi:G_E\to\operatorname{GL}_r(O_i)$ satisfying all
three fields of \texttt{Is\allowbreak Complement}.

Use the literal parameter list provided by the $Z_i$ membership
dictionary. It contains the distinguished part $(1,\mu,t)$.
Remove one occurrence of that part, retaining any other equal
occurrences, and let $T$ be the sum of the representations of
DFD5 for the remaining list. At a permitted Frobenius place the
original Hecke polynomial identity, the unramified sum law for
the isobaric parameter, and the coefficient-of-degree-$N-1$
calculation of DFD4 give
\[
 \iota_{K_i,\overline{\mathbb Q}_{\ell_i}}
            \operatorname{tr}R_\pi(g)
       =T(g)+\sum_{c=1}^{t}\chi_{i,c}(g).                \tag{DFD.9}
\]
Indeed the removed part contributes
$s_\mu(w)\sum_{k=0}^{t-1}q_w^{(t-N)/2-k}$, which is exactly
the displayed character sum. This verifies both its indexing
from $1$ to $t$ and its negative cyclotomic powers.

The left side of (DFD.9) is continuous under the supplied
$\ell_i$-adic coefficient embedding; the right side is
continuous by DFD5 and the given characters. DFD3 extends the
identity to all $g\in G_E$. Apply DFD1 and DFD2 with
$a=r$, $O=O_i$, and $K=K_i$. Their rational isomorphism is
precisely the source \texttt{Iso\allowbreak Block}; semisimplicity is that
of its complete block representation, and continuity is that
of the actual integral complement. Neither a trace-only
surrogate nor a representation over a different DVR is used.

The original \texttt{Complement\allowbreak F} and
\texttt{Complement\allowbreak Closed} follow with exactly their stated
quantifiers: apply this construction at each later member of
the given family, and then at every member of the given family
class. Their premises already supply \texttt{Galois\allowbreak Reps\allowbreak Cusp}
on that scope; no uniform choice of bases across distinct
coefficient rings is claimed or needed.
% END DFD6

% BEGIN DFD7
\paragraph{DFD7: one continuous point family, including points outside its support.}
\label{claim:DFD7}
Assume \texttt{Cheb}, \texttt{Galois\allowbreak Reps\allowbreak Cusp}, and
$\forall i\,\exists R_i\,\texttt{Good\allowbreak R}(D,i,R_i)$.
Choose these $R_i$ once. For each $\pi\in Z_i$ choose the
complement of DFD6 to $R_i(\pi)$. For $\pi\notin Z_i$ put
$\rho_i(\pi)(g)=1_r$ for every $g$.

Every $\rho_i(\pi)$ is continuous: this is DFD2 on the support
and continuity of a constant off the support. On $Z_i$ it is
the actual complement to the already chosen $R_i(\pi)$.
Thus the same two families $(R,\rho)$ satisfy the three
literal fields of \texttt{Pt\allowbreak Fam}:
\[
 \begin{gathered}
   \forall i\quad\texttt{Good\allowbreak R}(D,i,R_i),\\
   \forall i,\pi\quad\rho_i(\pi)\text{ has continuous matrix entries},\\
   \forall i,\pi\in Z_i\quad
       \texttt{Is\allowbreak Complement}(R_i(\pi),\rho_i(\pi)).
 \end{gathered}                                         \tag{DFD.10}
\]
This proves the whole original
\texttt{exists\_\allowbreak continuous\_\allowbreak complement\allowbreak Family} and
\texttt{exists\_\allowbreak pt\allowbreak Fam}. Classical choice is only over the
specified location and point sets. Continuity here is in the
Galois variable, so choosing the families does not create an
unproved assertion of continuous dependence on $i$ or $\pi$.
% END DFD7

\subsection{The literal canonical data record}

% BEGIN DFD8
\paragraph{DFD8: all nine canonical row-data laws.}
\label{claim:DFD8}
Fix $L:\texttt{Size\allowbreak Data}_2(D)$, a supplied
$HW:\texttt{HWOf}(D.\mathrm{loc})$, the row-constant record
$R=(C_2,C_B,C_C,C_A,C'_B,C_K)$, and a \texttt{Pt\allowbreak Fam}
$(R_i,\rho_i)$. This paragraph's letter $R_i$ denotes the
representation; $R$ alone denotes the row-constant record.
Put
\[
 \begin{split}
 K_1&=C_B+48,\quad K_2=6C_B+12+6C_C,\\
 K_3&=3C_AC_B+12C_A+9,\qquad h=\lfloor2b\rfloor+1,\\
 c_B(i)&=\texttt{c\allowbreak Bnum}(\ell_i,\lfloor2b\rfloor,t,4nn_V),\\
 c_0(i,v)&=5+\lfloor\log_{\ell_i}(2h)\rfloor
       +4nn_V\bigl(4+2\lfloor\log_{\ell_i}(2h)\rfloor\bigr).
 \end{split}                                           \tag{DFD.11}
\]
In the last line the floors denote the source natural-number
logarithm, so the formula has its prescribed integer meaning.
Choose $\epsilon_J$ to be the witness of
\texttt{Conclusion\allowbreak Ja} when that proposition holds, and choose
the Boolean value true otherwise. Define $dc_S(i)$ to be a
witness of $\texttt{hyp\allowbreak Decl}(i,dc,\epsilon_J)$ when both
\texttt{Split\allowbreak Row\allowbreak Prime} and such a witness hold; otherwise take
all places over $\ell_i$ as rows and no columns. This is exactly
$\texttt{dc\allowbreak Of}$ and is defined without assuming these
propositions true.

For arbitrary $e'(i,v)$, finite $\mathrm{exc}\subset\mathbb N$,
and $I:\texttt{Tw\allowbreak In}(D)$ construct $X$ as follows. Keep
$X.r=(R_i)_i$, $X.\rho=\rho$, $X.e'=e'$, $X.c_B=c_B$,
$X.c_0=c_0$, $X.\mathrm{exc}=\mathrm{exc}$,
$X.\eta=\lfloor L.\mathrm{etaSz}\rfloor$, and $X.I=I$.
The band is the actual predicate
$\texttt{band\allowbreak Of}(L,HW,C_B,i,R_i,\rho_i,v)$, and
$X.e_S(i)=\texttt{e\allowbreak A}(L,C_A,\ell_i)$. Write
\[
 \begin{split}
 B_i&:\quad\ell_i<\texttt{l\allowbreak A}(L)
                    \ \text{and}\ 4L_w<\texttt{ord\allowbreak H}(D,\ell_i),\\
 U_i&=\begin{cases}
 \texttt{len\allowbreak Row\allowbreak B}(L,C_2,\ell_i),&B_i,\\
 \texttt{len\allowbreak Row\allowbreak CD}(L,X.\mathrm{band}_i,K_1,K_2,K_3,\ell_i),&\neg B_i.
 \end{cases}
 \end{split}                                           \tag{DFD.12}
\]
Set $X.\mathrm{lenRow}_i=U_i$. If $\ell_i$ is inert in the
source sense, use the all-row declaration, put
$X.\mathrm{len0}(i,\pi,v)=U_i(v)$, and use
$\texttt{local\allowbreak Error\allowbreak Bound\allowbreak Decomposition}(i,X.dc_i,c_B(i))$
for the error. Otherwise use $X.dc_i=dc_S(i)$,
$\texttt{len\allowbreak A}(i,X.dc_i,X.e_S(i),v)$ for $\mathrm{len0}$,
and the actual
$\texttt{local\allowbreak Error\allowbreak Bound\allowbreak Places}(i,\rho_i,\pi,X.dc_i,c_0(i),c_B(i))$
for the error. Rational row lengths are cast to real numbers
only in the two real-valued fields.

These formulas verify each of the nine \texttt{Data\allowbreak AD} laws.
The inert declaration has empty column set. Its $B_i$ and
$\neg B_i$ branches give \texttt{row\allowbreak B} and \texttt{row\allowbreak CD}.
The inert length and error definitions give \texttt{len0I}
and \texttt{error\_\allowbreak bound\_\allowbreak inert\_\allowbreak eq}. The formula for $e_S$
gives \texttt{e\allowbreak S}. A split prime cannot be inert: the
original prime-class definition assigns disjoint classes to
these predicates. Hence the noninert length and error branches
give \texttt{len0S} and \texttt{error\_\allowbreak bound\_\allowbreak split\_\allowbreak eq}.
The actual band formula is the ninth law. On every noninert
location $X.dc=dc_S$, as required separately by
\texttt{Canon}. Together with (DFD.10) and the three exact
equalities for $c_B,c_0,\eta$, this proves all six fields of
the original \texttt{Canon}. No row bound or crystalline
comparison is inferred merely from these definitions.
% END DFD8

% BEGIN DFD9
\paragraph{DFD9: the whole original canonical-data existence theorem.}
\label{claim:DFD9}
Assume \texttt{Cheb}, $\texttt{adm2}(D)$, and
$\texttt{Conclusion41}(D)$. For every original $L,HW,R$,
every function
\[
 e':\texttt{RFam}(D)\to\texttt{Rho\allowbreak Fam}(D)\to
          \prod_i\bigl(\operatorname{Places}(E)\to\mathbb N\bigr),
 \quad
 \mathrm{exc}:\texttt{Rho\allowbreak Fam}(D)\to\operatorname{Finset}(\mathbb N),
                                                               \tag{DFD.13}
\]
and every $I:\texttt{Tw\allowbreak In}(D)$, there is an actual $X$ with
\[
 \begin{gathered}
 \texttt{Canon}(L,HW,R,X),\qquad X.e'=e'(X.r,X.\rho),\\
 X.\mathrm{exc}=\mathrm{exc}(X.\rho),\qquad X.I=I,
 \qquad\forall i\ \texttt{Good\allowbreak R}(D,i,X.r_i),\\
 0\le L.\mathrm{etaSz}\ \Longrightarrow\
          X.\eta\le2L.\mathrm{etaSz},\\
 \forall i\quad
 \texttt{c\allowbreak Bnum}(\ell_i,\lfloor2b\rfloor,t,4nn_V)\le X.c_B(i).
 \end{gathered}                                         \tag{DFD.14}
\]
DFD4 supplies \texttt{Galois\allowbreak Reps\allowbreak Cusp}. For every $i$, the
first existence clause of \texttt{Conclusion41} supplies
some $R_i$ with \texttt{Good\allowbreak R}. DFD7 chooses one pair of
families $(R,\rho)$. Only after that choice evaluate the
functions in (DFD.13), obtaining $e'(R,\rho)$ and
$\mathrm{exc}(\rho)$. Apply DFD8 with these values.
Its formulas prove the first four assertions in (DFD.14),
and its \texttt{Pt\allowbreak Fam.good} proves the fifth. For the sixth,
$\lfloor a\rfloor_{\mathbb N}\le a\le2a$ when $a\ge0$.
The final inequality is equality by (DFD.11).

This order of construction preserves the original functional
dependence of the error and exceptional set on the chosen
representations. No prime condition is added to the finite
set of natural numbers. This is the full source theorem
$\texttt{Global\allowbreak Theory.Hub\allowbreak X.exists\_\allowbreak hub\allowbreak Data}$, with every
original input and output clause.
% END DFD9

\subsection{The same rows, twists, and coefficient lattices}

% BEGIN DFD10
\paragraph{DFD10: an actual good row for each good point.}
\label{claim:DFD10}
The original \texttt{tie0}, \texttt{Conclusion42}, and
\texttt{Conclusion\allowbreak Z} imply \texttt{Good\allowbreak Pt\allowbreak Row}. In fact
only the last input and the built-in $D$ dictionaries are
needed for this implication.

For $\pi\in Z_i^{\mathrm{good}}$, the literal good-point
dictionary supplies a parameter list $\Sigma$ with
$\texttt{Has\allowbreak Param\allowbreak X}(i,\pi,\Sigma)$ and
$\texttt{good\allowbreak Sig}(\Sigma)$. Its actual Hecke eigenspace
$\mathrm{eigC}(i,\pi)$ is nonzero: the witnessing
automorphic representation has nonzero type space, and that
type space lies in the supremum defining this eigenspace.
The \texttt{Conclusion\allowbreak Z} equality identifies this eigenspace
with the image of
\[
 M_{i,\Sigma}
    =\bigvee_{\sigma:\,\texttt{Has\allowbreak Param\allowbreak W}(i,\sigma,\Sigma)}
                      \mathrm{isot}(j_i,\sigma).         \tag{DFD.15}
\]
Therefore $M_{i,\Sigma}\ne0$. If every displayed isotypic
space were zero, its supremum would be zero. Choose one
$\sigma$ with \texttt{Has\allowbreak Param\allowbreak W} and nonzero isotypic space.

At each designated type place $w_{k,v}$ the original
nonconjugacy condition implies it is split. Explicitly, if
conjugation carried its prime ideal into itself, maximality
would make the ideal equal to its conjugate, contrary to
nonconjugacy; this is the split predicate. The split local
identity in \texttt{Has\allowbreak Param\allowbreak W} consequently identifies the
label of $\sigma$ at every such place with the label computed
from $\Sigma$. Unfolding the two goodness predicates now
gives
$\texttt{good\allowbreak Row}(\sigma)\Longleftrightarrow
\texttt{good\allowbreak Sig}(\Sigma)$, one label condition at a time.
Thus the chosen $\sigma$ is good. Together with its nonzero
isotypic space and the same $\Sigma$, this proves every field
of \texttt{Good\allowbreak Pt\allowbreak Row}. If \texttt{Triples\allowbreak Good} is also
supplied, apply it to this very pair $(\sigma,\Sigma)$ to
obtain a good signature attached to $\pi$, with all the
original shared parameters unchanged.
% END DFD10

% BEGIN DFD11
\paragraph{DFD11: the exact chosen-twist killing implication.}
\label{claim:DFD11}
For the original $D$, $C:\texttt{NChoice}(D)$,
$I:\texttt{Tw\allowbreak In}(D)$, and $\rho:\texttt{Rho\allowbreak Fam}(D)$,
\[
 \texttt{NCh\allowbreak Tie}(D,C)\ \wedge\ \texttt{Tw\allowbreak Good\allowbreak At}(D,C,I)
 \ \wedge\ \texttt{Signature\allowbreak Choice\allowbreak Condition}(\rho,C)
       \ \Longrightarrow\ \texttt{Kill\allowbreak At}(\rho).        \tag{DFD.16}
\]
Fix a twisted location $i$ and put $\ell=\ell_i$.
The tied twist index has prime $\ell$, and its error budget
is the original $\texttt{tw\allowbreak Eps}=\texttt{eps\allowbreak L}(\ell)=1/\ell^2$.
The \texttt{Tw\allowbreak Good\allowbreak At} data give nonnegative real numbers
$\epsilon_s$ on the finite good-signature set, with
\[
 \sum_{s\text{ good}}\epsilon_s\le1/\ell^2,
 \qquad
 \mathrm{sLen}(s)\le
          \operatorname{ofReal}(\nu+\epsilon_s),         \tag{DFD.17}
\]
using the exact oriented character in the source. More
explicitly, put
\[
 \psi_W=(\texttt{psi\allowbreak At}(C.\mathrm{sd},C.\psi'))
                       \circ\mathrm{art}.
\]
Both the bound in \texttt{GoodTw.partAT} and the bound in
\texttt{SignatureChoiceCondition} use
$\texttt{psi\allowbreak At}(\mathrm{side}(s),\psi_W)$.
They therefore compare the identical character directly.
The two Boolean sides need not agree, and no cancellation
of these orientations is used.

For $\pi\in Z_i^{\mathrm{good}}$, the supplied signature
condition chooses a good $s$ with
\[
 \operatorname{ofReal}(\mathrm{sPi}(i,\pi))
    \le\mathrm{sLen}(s),
 \qquad\mathrm{sLen}(s)\ne\infty
                \Longrightarrow\mathrm{sFin}(i,\pi).
\]
The right side of (DFD.17) is finite, so $\mathrm{sFin}$
holds. Every $\epsilon_s$ is bounded by the sum in
(DFD.17), because all summands are nonnegative. Since
$\nu+\epsilon_s\ge0$, comparison of the finite extended
nonnegative reals gives
\[
       \mathrm{sPi}(i,\pi)
          \le\nu+\epsilon_s\le\nu_L(D,\ell)+1/\ell^2.     \tag{DFD.18}
\]
Here $\nu=\nu_L(D,\ell)$ is the exact host cyclotomic
identity in the original \texttt{Nu\allowbreak Host} definitions, not
the weaker logarithmic bound used for DNE2. Explicitly,
unfold \texttt{nu\allowbreak L}, select its prime branch using the
primality of the tied twist prime, and the remaining
expression is definitionally \texttt{Gen.nu}. Thus this
step uses no separate numerical estimate.

For the second field put
\[
 d_Q=\sum_{v\in Q_i}
       \mathrm{lenUnit}_{O_i,\ell}
                       (O_i/(\zeta_{i,v}-1)).
\]
The signature condition gives
$d_Q\le\sum_{P\in C.Q(w_i)}\mathrm{tameVal}(\psi_W,P)$.
The last \texttt{Tw\allowbreak Good\allowbreak At} inequality bounds $t$ times
this same tame sum by $1/\ell^2$. Multiplication by the
nonnegative integer $t$ gives $td_Q\le1/\ell^2$, exactly
the second field of \texttt{Kill\allowbreak At}. This proves the full
original \texttt{twist\_\allowbreak kills\_\allowbreak selmer} theorem. Producing
\texttt{Signature\allowbreak Choice\allowbreak Condition} for the chosen $\rho$
is a separate comparison problem and is not inferred from
the existence of an attached signature alone.
% END DFD11

% BEGIN DFD12
\paragraph{DFD12: finite coefficients from the original integral crystalline input.}
\label{claim:DFD12}
The source predicate \texttt{Bhatt\allowbreak Scholze\allowbreak Integral} implies
the complete predicate \texttt{Bhatt\allowbreak Scholze\allowbreak Crystalline\allowbreak Lattices}.
The retained input is the former predicate, including its
fully faithful assertion; no integral comparison theorem
is newly assumed in the conclusion.

Fix the original prime $\ell$, local base field, prism,
$A_{\inf}$ identification and standard period data. Let
$L/\mathbb Q_\ell$ be any finite coefficient extension and
$T\subset V$ an $\mathcal O_L$-lattice in the given
continuous crystalline $L$-representation. Restrict the
same subset $T$ to a $\mathbb Z_\ell$-module. The ring
$\mathcal O_L$ is finite over $\mathbb Z_\ell$, so $T$ is
finite over $\mathbb Z_\ell$. It is torsion-free, as it
lies in a $\mathbb Q_\ell$-space, and hence free over the
PID $\mathbb Z_\ell$.

It spans $V$ over $\mathbb Q_\ell$. To see the point
without changing the lattice, every $c\in L$ has the form
$c=q a$, $q\in\mathbb Q_\ell$ and $a\in\mathcal O_L$:
multiply by a sufficiently large power of $\ell$ to make
it integral. Therefore the $\mathbb Q_\ell$-span of $T$
is stable under $L$. Since its $L$-span is $V$, this
$\mathbb Q_\ell$-span is already $V$. Stability under the
Galois action is unchanged. The finite-dimensional
$L$-module topology is the same module topology after
restriction to $\mathbb Q_\ell$, so every original orbit
map remains continuous.

For completeness, crystalline admissibility also survives
this restriction. In the source period-space definition,
$D_{\mathrm{cris}}$ is cut out by the same equations
$(g-1)x=0$ before and after restriction. Thus
\[
 D_{\mathrm{cris}}(\operatorname{Res}V)
   =\operatorname{Res}D_{\mathrm{cris}}(V).              \tag{DFD.19}
\]
Each dimension over $\mathbb Q_\ell$ is
$[L:\mathbb Q_\ell]$ times its dimension over $L$.
Substitute this on both sides of the source crystalline
dimension equality (including its unchanged base-period
degree), then cancel the positive factor
$[L:\mathbb Q_\ell]$. This proves the exact
$\texttt{is\allowbreak Crys\_\allowbreak res\allowbreak Q\_\allowbreak iff}$, rather than assuming it.

Apply \texttt{Bhatt\allowbreak Scholze\allowbreak Integral} to this restricted
representation and this identical stable lattice. It
returns the required $F$-crystal, its semilinear realization
map, and its realization law with target
$W(C^\flat)\otimes_{\mathbb Z_\ell}T$. These are already
the output fields of the coefficient version. Its second,
unique-morphism field is literally the second field of
the retained integral predicate, so it passes through
unchanged. This completes the finite-coefficient theorem.
Constructing the compatible Wach functor over all layers
and its \texttt{Cmp\allowbreak Of} ties requires further work; it is
not a consequence of scalar restriction alone.
% END DFD12

\subsection{One eventual threshold and its exact remaining inputs}

% BEGIN DFD13
\paragraph{DFD13: all original parameter inequalities at one uniform cutoff.}
\label{claim:DFD13}
For every $\epsilon>0$ and real $A$ there is $x_0$ such
that every original $(D,L,P)$ with scale exponent
$\epsilon$, $P.C_{L_w}\le A$, and $x\ge x_0$ satisfies
both \texttt{Adm\allowbreak Par} and
\texttt{Parameters\allowbreak Polynomially\allowbreak Bounded}. All objects
$(D,L,P)$ are quantified after this one cutoff.

Write $b=e^{x^\epsilon}$,
$b'=xb/e^{x^{\epsilon/3}}$,
$w''=b'/(xn)$, $S=L.\mathrm{Sc}$, and
$\eta=L.\mathrm{etaSz}$. The original parameter fields
give $t\le\eta$, $t\ge7$, $n\ge1$,
$kS\le w''$, $xt\eta\le w''$ and
$k\ge m\sqrt{x}$ with $m\ge2$.
For $x\ge10000$, these imply
\[
 200S\le w'',\qquad70000\eta\le w'',\qquad
 100\le w'',\qquad100(\eta+S+t+2)\le w'',
 \qquad w''\le b'.                                    \tag{DFD.20}
\]
For the fourth inequality, $t\le\eta$ and $\eta\ge7$
give $\eta+t+2\le(16/7)\eta$; its contribution after
multiplication by $100$ is less than $w''/2$ by the
second inequality, while $100S\le w''/2$.

Here is also the top-weight inequality, with its actual
parameter inputs. Put
$T_0=\texttt{Nstar}(L)x^\epsilon\mathrm{LaG}/k$.
The field \texttt{rho\_\allowbreak le}, after dropping its other
nonnegative term, says $4T_0\le x^{-2}$. Nonnegativity
of that term follows from $\mathrm{LaG}>0$, which itself
follows from $4S\le\mathrm{LaG}\,b'/k$ and $S>0$.
The field \texttt{b\allowbreak T\_\allowbreak le} gives $b_a\le b'(1+T_0)$.
For $x\ge4$,
\[
        b_a+b'/x\le b'(1+1/64+1/4)<3b'/2.              \tag{DFD.21}
\]
Finally, uniformly in $P.C_{L_w}\le A$,
$L_w\le\max(A,0)\log x$. The elementary limit
\[
 x(\max(A,0)\log x+1)e^{-x^{\epsilon/3}}\longrightarrow0
\]
provides a cutoff for
\[
              100(L_w+1)b'\le b.                       \tag{DFD.22}
\]
To justify that limit directly, put $y=x^{\epsilon/3}$.
The expression is bounded for $x\ge1$ by a constant
times $y^{6/\epsilon}e^{-y}$; expand $e^{y/2}$ in its
power series to dominate any fixed power of $y$, leaving
a constant times $e^{-y/2}\to0$.

Equations (DFD.20)--(DFD.22) are all fields of the source
$\texttt{Roomy}(L,100)$. They give
$\eta+t+S<w''$, $3w''\le b$, every
$b_a+b'/x\le b$, and $t\le\eta$, exactly
\texttt{Adm\allowbreak Par}. They also give $S+\eta+t+2\le b$,
exactly \texttt{Parameters\allowbreak Polynomially\allowbreak Bounded}. Taking
the maximum of the numerical and exponential cutoffs
proves the full original
\texttt{admissible\allowbreak Parameters\_\allowbreak of\_\allowbreak large\_\allowbreak scale} theorem.
% END DFD13

% BEGIN DFD14
\paragraph{DFD14: the full family small-prime assertion.}
\label{claim:DFD14}
For an original family $F$ and $C\ge1$, there is one
$x_0$ such that, for every $x\ge x_0$, every original
$h_x:F.x_1<x$, and every $i\in F.I_x$, one has
\[
 \texttt{Small}(L,C),\qquad8\le L.\mathrm{sc}.L,
 \qquad60\le D.\mathrm{host}.\mathrm{splitBound},
 \qquad C\le b.                                        \tag{DFD.23}
\]
The proof below verifies all fields of \texttt{Small},
including the prime-three and two auxiliary-prime clauses.

Set $z=x^\epsilon$, $b=e^z$, $p_*=\texttt{l\allowbreak Star}(L)$,
$A_*=\texttt{l\allowbreak A}(L)$, and use the original rank symbols
$f,r,m,t$. The family size bounds give, for $x\ge3$,
\[
 k\le x^5,\quad f\le k,\quad r\le2k,
 \quad t\le x^{|C_t|+1},\quad
 frt\le2x^{11+|C_t|}.                                  \tag{DFD.24}
\]
The exponent $C_t$ is the one fixed by the family, before
its members. Choose a cutoff, depending only on the
family constants and $C$, such that
\[
 \begin{gathered}
 z\ge\max(16,C),\qquad x^{\epsilon_z/5}\ge60,\\
 4\cdot2^C x^{(11+|C_t|)C+\epsilon}\le e^{x^\epsilon},
 \qquad32x^{5+\epsilon}\le e^{x^\epsilon}.
 \end{gathered}                                         \tag{DFD.25}
\]
Every one of these conditions is eventually true since
$\epsilon,\epsilon_z>0$. For the exponential comparisons
use the same power-series argument as DFD13. They are
uniform in every family member.

The original scale and parameter definitions give
\[
 A_*\ge z,\quad S,\eta,t\le b',\quad
             rb'A_*\log x\le mb.                       \tag{DFD.26}
\]
For clarity these do not require a new asymptotic result:
$A_*=kz/(m\delta)$, $k\ge m$, $0<\delta\le1$ give the
first; $t\eta x\le w''\le b'$, $t\le\eta$, and
$kS\le w''$ give the middle bounds for $x\ge1$; and
the field \texttt{rb\allowbreak S}, after dropping its nonnegative
$S$ term and using $\delta\le1$, gives the last.
Since $\log x\ge1$ for $x\ge3$, (DFD.25)--(DFD.26)
give
\[
 16rb'\le mb,\quad16rz\le mb,\quad
 r(S+t+\eta+z)\le mb/4,\quad(frt)^Cz\le b/4.            \tag{DFD.27}
\]
The second follows from $r\le2x^5$ and the final
inequality in (DFD.25); the first uses $A_*\ge16$.

For either auxiliary prime $p=q,d$, the host inequality
$p>4tr+1$ gives $2r\le p$. Consequently $rb/p\le b/2$.
Using $f,m\ge1$, the actual expression
\[
 \texttt{mid\allowbreak G}(L,C,p)
   =Cf\{r(b/p+S+t+\eta+z)+mb\}+C(frt)^Cz
\]
is at most
$Cf(7mb/4)+C(b/4)\le2Cfmb$, by (DFD.27).
This is precisely the auxiliary-prime field of
\texttt{Small}.

It remains to locate the exceptional small primes. From
\texttt{lstar1} one has $b/p_*<p_*$, hence
$\sqrt b\le p_*$. The original conductor condition is
$M^{36}t^{150}\le\sqrt b$, where $M=Q_Sqd\ge1$.
Every prime in $D.\mathrm{host}.\mathrm{small}$ and
both $q,d$ divide the positive integer $M$. Thus each
is at most $M\le p_*$. The original \texttt{lstar2}
also gives $7<p_*$, in particular $3\le p_*$.

The split threshold in (DFD.25), together with the actual
floor tie for \texttt{split\allowbreak Bound}, gives
$60\le\texttt{split\allowbreak Bound}$. Hence the original host
small-splitting law applies to $3$. This prime differs
from $q$ by the host record and from $d$ because
$d>4tr+1>21$. If it belongs to the host's small set, its
assigned class is noninert; otherwise the just obtained
split law assigns its noninert split class. Thus
$\neg\texttt{Inert\allowbreak K}(D,3)$. We have proved every small,
auxiliary, prime-three, and \texttt{mid\allowbreak G} clause.
Finally $L.\mathrm{sc}.L=z\ge16$ and
$b=e^z\ge z\ge C$. This proves (DFD.23) and the full
original \texttt{special\_\allowbreak primes\_\allowbreak small} family theorem.
% END DFD14

% BEGIN DFD15
\paragraph{DFD15: the original family setup reduced to two actual suppliers.}
\label{claim:DFD15}
Fix the original $F:\texttt{Fam\allowbreak P}(\epsilon)$, constants
$C,C_2,CT$, and the original \texttt{Tw\allowbreak Good} and
\texttt{Facts\allowbreak F} premises. Retain \texttt{Cheb} and the
following two explicit supplier statements on these same
data:
\begin{enumerate}
\item For each relevant $D$ there exist actual $HW,F_W$
and $\texttt{Cmp\allowbreak Of}(D.\mathrm{loc},D.\mathrm{locA},HW,F_W)$.
\item For each actual pair $(R,\rho)$ satisfying
\texttt{Pt\allowbreak Fam} on this $D$, the original facts imply
$\texttt{Signature\allowbreak Choice\allowbreak Condition}(\rho,C_D)$ for the
same prescribed $C_D=\texttt{nch}$.
\end{enumerate}
Then the full output of \texttt{exists\_\allowbreak setup\_\allowbreak large}
follows, with one $x_0$ before $x,h_x,i$, before the
$\texttt{Sha\allowbreak Fin\allowbreak Odd}$ proof, and before the arbitrary
finite exceptional set. This is a conditional reduction;
the two displayed statements are remaining human-proof
obligations, not additional consequences claimed here.

The nonempty index at scale zero shows $\epsilon>0$
from the family's scale tie. DFD13 gives a cutoff $x_a$
with $A=F.PK[7]$, because that is exactly the family
bound for $C_{L_w}$. DFD14 at constant $1$ gives $x_b$.
The original \texttt{Tw\allowbreak Good} premise gives $x_c$.
Choose $x_0=\max\{x_a,x_b,x_c\}$ once.

For a later $x,h_x,i$, let $D,L,C_D$ be its original
data. The twist premise chooses $N_0,S$ and gives
\texttt{Tw\allowbreak Good\allowbreak At} for
$I=(\texttt{tw\allowbreak E}(F,x,h_x,i),N_0,S)$. Keep this $I$.
The first remaining supplier chooses $HW,F_W,\mathrm{Cmp}$.
The \texttt{Facts\allowbreak F} record contains \texttt{adm2} and
\texttt{Conclusion41}, so DFD9 constructs $X$ with
row constants
\[
            \texttt{row\allowbreak C}(C_2)=(C_2,320,25,23,32,32),
\]
$e'=0$, exceptional-set function constantly equal to
the given $\mathrm{exc}$, and exactly this $I$.
Thus $X.\mathrm{exc}=\mathrm{exc}$ and
$X.I=I$. The second supplier applies to its actual
\texttt{Pt\allowbreak Fam}; DFD11 and the original \texttt{NCh\allowbreak Tie}
give $\texttt{Kill\allowbreak At}(X.\rho)$.

Now assemble the eight \texttt{Setup} fields.
The facts field is the given \texttt{Facts\allowbreak F} member;
the canonical field is DFD9; its two parameter fields
are DFD13; its $8\le L$ and split-bound fields are
DFD14; its killing field was just proved; and its last
field is the supplied \texttt{Sha\allowbreak Fin\allowbreak Odd}. The
\texttt{Adm\allowbreak J} conclusion follows from the same member's
original \texttt{adm\allowbreak J} implication applied to
\texttt{Adm\allowbreak Par}. All fields refer to the one chosen
$X,HW,F_W,\mathrm{Cmp},D,L,C_D$.

This proves the reduction with the exact existential and
universal order. DFD12 supplies the finite-coefficient
layer of the first remaining supplier from the original
integral input; DFD10 supplies actual attached good rows
for the second. Neither yet supplies all their coherent
comparison laws. Later local row estimates and the full
common \texttt{Concl\allowbreak F} also remain. Accordingly this
chapter does not close \texttt{GLOBAL\_\allowbreak DENOMINATOR}.
% END DFD15

\subsection{Source alignment and remaining scope}

The exact pinned source ranges and full-file hashes for
each DFD claim are in this component's
\texttt{source\_\allowbreak manifest.json} and \texttt{provenance.json}.
DFD1--DFD7 follow the actual interfaces in
\texttt{Semisimple\allowbreak Of\allowbreak Injective},
\texttt{Block\allowbreak Complement\allowbreak Lattice},
\texttt{Galois\allowbreak Reps\allowbreak Cuspidal}, and
\texttt{Exists\allowbreak Continuous\allowbreak Complement\allowbreak Family}; the saturated
lattice proof supplies a direct simplification of the
matrix construction. DFD8--DFD9 prove the literal
\texttt{Canon} and \texttt{Exists\allowbreak Hub\allowbreak Data} output.
DFD10--DFD12 correspond to \texttt{Local\allowbreak Global\allowbreak Assembly},
\texttt{Chosen\allowbreak Twist\allowbreak Annihilates\allowbreak Selmer\allowbreak Groups}, and
\texttt{Bhatt\allowbreak Scholze\allowbreak Coefficients\allowbreak Of\allowbreak Padic\allowbreak Int\allowbreak Lattices}.
DFD13--DFD14 supply the original parameter and small-prime
theorems. DFD15 records precisely what is still needed
to replace the conditional use of
\texttt{Exists\allowbreak Hub\allowbreak Data\allowbreak Family} by a complete human proof.

The earlier DNE chapter remains the full
$\texttt{Concl\allowbreak F}\Rightarrow\texttt{Adm\allowbreak F}\Rightarrow
\texttt{Link\allowbreak N}$ consumer. The present actual integral
families feed its upstream construction without changing
the host, coefficient DVRs, embeddings, chosen twists,
or exceptional sets. No formal replay is claimed.
